\par\SourceLeaf{144}
A subtle scrutiny of the wording might perhaps find reasons in favour of the first hypothesis; but I shall dwell on this all the less, since I myself cannot regard that hypothesis as admissible, chiefly because P. would otherwise contradict his own principles (above, p. [133]). I believe rather that, misled by \textsc{Kritter}'s opinion, — or, as one may also say, and as I shall soon show in detail, by his mistaken understanding of this opinion — he regarded the number of 15 pensions as an impassable or almost impassable barrier, under whose protection he could make his plan in complete safety.

I find a reflection of similar confidence in E.'s vote, which is distinguished by being the only one to mention the fatalistic number 15, and whose essential content, insofar as it belongs here, is as follows.

F.'s two questions 1 and 2, says E., must be decided solely by the actual fund and its yield, compared with \textit{the tested} principles so convincingly presented by H. P. It might be that the extended duration of the orphans' pensions would cause the highest number of widows, 15, to be exceeded, which, as incompatible with the fund's safety, must absolutely never be allowed to happen. He therefore considers it advisable to approve the extension of the orphans' pensions only with the express restriction, \textit{as long as the number of pensions is not thereby raised above 15}. This restriction seems all the more necessary because the calculation of the children's circumstances is not supported by experience \textit{as certain} as the observed proportion of widows. I understand this as follows: E., whose understanding of P.'s plan is evidently the exact opposite of \textsc{Sch.}'s, regards it as certain, according to the \textsc{Kritter}-P. theory, that the number of widows can never exceed the calculated number (thus 13); but it is not equally certain that no more than 2 orphans' pensions could be added, and in case this did happen and the total number of pensions would thereby be driven above 15, the restriction of orphans' pensions would have to be reserved.

Similar fears that the number of orphans' pensions might have been estimated too low were also expressed by some other voters; but they were neither acted upon in the final decision nor confirmed by experience so far. In fact, from the establishment of the widows' fund until now, there have never yet been more than two orphans' pensions at the same time. I do not, however, wish to leave unmentioned that during a short part of the year 1776 the number would have risen to 3 if the extension of orphans' pensions until the completion of the \(20^{\text{th}}\) year had already been in force then.

With the \textit{number of widows}, however, matters stand quite differently. The supposed certainty that it could not rise above 13 has been destroyed by more recent experience, since it has already once risen to 22, and even its average during the last 8 years has been somewhat above 19. Long before 1794 the number of widows (again, that is, without counting the orphans' pensions) had once risen to 14, and during a period of five years (1774—1779) had maintained an average of 13. This fact, which could have been known to the curator but certainly not to the voting professors, could, if considered from the only admissible standpoint indicated above on p. [130], l. [4], already have served then to prove the incorrectness of P. and E.'s assumption.

But since the assumption of the number 13 as the highest number of widows rests on \textsc{Kritter}'s opinion, in which half the number of existing marriages is put forward as determining the maximum number of widows, the conclusion seems natural that there may be errors in it, and I have therefore reached the point at which I must subject this opinion itself to criticism.
\par\SourceLeaf{145}\noindent
To make this criticism completely intelligible, I am compelled to give some developments beforehand, which some readers might perhaps find it laborious to follow me through if it is not yet apparent from the outset what they will lead to. It will therefore, it seems to me, be appropriate to put the principal objections to \textsc{Kritter}'s opinion and the conclusions drawn from it at the very beginning here.
\Needspace{28\baselineskip}
\begin{list}{}{\setlength{\leftmargin}{2em}\setlength{\labelwidth}{1.5em}\setlength{\labelsep}{.5em}\setlength{\itemsep}{0pt}\setlength{\parsep}{0pt}}
\item[1)] The applicability of the ratio of 1 widow to 2 existing marriages to the professors' widows' fund here is not certain; rather, as already remarked above on p. [133], l. [5], it is probable that the circumstances of this fund require one to reckon on somewhat more widows.
\item[2)] \textsc{Kritter}'s assertion [l. 10 of the reprint], even when put forward independently of this assumed ratio \(2:1\), namely that one may assume an average of 10 years already sufficient in every case to indicate very nearly the true ratio applicable to the fund's circumstances, is entirely false, and the incorrectness of this proposition can be demonstrated on theoretical grounds alone, even without appealing to experience.
\item[3)] Far more important than these two objections is the fact that P. has misinterpreted \textsc{Kritter}'s opinion, since what P. understood by the maximum number of widows, and what is denoted by this word when a definite standard ratio between existing marriages and widows is established, and what is actually meant in \textsc{Kritter}'s opinion,
\begin{center}\textit{are two very different things.}\end{center}
To this is added the equally important circumstance
\item[4)] that the numerical results which were admissible for a society of a certain size remaining nearly the same throughout must be substantially altered when that size expands considerably, as has already been discussed in detail above.
\end{list}
\par\medskip\centerline{\rule{.25\linewidth}{.3pt}}\medskip
